\chapter{KATA PENGANTAR}
		\begin{figure}[h]
			\centering
			\includegraphics[width=0.5\linewidth]{img/bismillah.png}
		\end{figure}

		Alhamdulillahirabbil’alamin, segala puji bagi Allah SWT, yang telah melimpahkan rahmat dan hidayah-Nya sehingga penulis dapat menyelesaikan Tugas Akhir yang berjudul \textbf{Rancang Bangun \textit{Middleware} untuk Melacak Konfigurasi Perangkat Jaringan Menggunakan Git}. Pengerjaan Tugas Akhir ini merupakan suatu kesempatan yang sangat baik bagi penulis. Dengan pengerjaan Tugas Akhir ini, penulis bisa belajar lebih banyak untuk memperdalam dan meningkatkan apa yang telah didapatkan penulis selama menempuh perkuliahan di Departemen Informatika ITS. Dengan Tugas Akhir ini penulis juga dapat menghasilkan suatu implementasi dari apa yang telah penulis pelajari.
		Selesainya Tugas Akhir ini tidak lepas dari bantuan dan dukungan beberapa pihak. Sehingga pada kesempatan ini penulis mengucapkan syukur dan terima kasih kepada:
		\begin{enumerate}
			\item Allah SWT atas anugerahnya yang tidak terkira kepada penulis dan Nabi Muhammad SAW.
            \item Keluarga penulis yang selalu menyemangati.
			\item Royyana Muslim Ijtihadie, S.Kom., M.Kom.,   Ph.D selaku pembimbing I yang telah membantu, membimbing dan memotivasi penulis mulai dari pengerjaan proposal hingga terselesaikannya Tugas Akhir ini.
			\item Bapak Bagus Jati Santoso, S.Kom., Ph.D selaku pembimbing II yang juga telah membantu, membimbing dan memotivasi penulis mulai dari pengerjaan proposal hingga terselesaikannya Tugas Akhir ini.
            \item Teman-teman \textit{Administrator} laboratorium AJK.
			\item Darlis Herumurti, S.Kom., M.Kom., selaku Kepala Departemen Informatika ITS pada masa pengerjaan Tugas Akhir, Bapak Radityo Anggoro, S.Kom., M.Sc., selaku koordinator TA dan segenap dosen Departemen Informatika yang telah memberikan ilmu dan pengalamannya.
			\item Serta semua pihak yang telah turut membantu penulis dalam menyelesaikan Tugas Akhir ini.
		\end{enumerate}

		Penulis menyadari bahwa Tugas Akhir ini masih memiliki banyak kekurangan. Sehingga dengan kerendahan hati, penulis mengharapkan kritik dan saran dari pembaca untuk perbaikan ke depannya.

	\hfill Surabaya, Januari 2020 \\ \\ 


  \hfill Muhammad Faris Didin A.